% !TEX program = xelatex
% 俄语 · 名词总表（按课件分类）· 2026-11-13
\documentclass[11pt]{ctexart}
%=====================================================================
% theme-violet.tex —— 讲义精讲 · 紫花主题
% 用法：\documentclass[11pt]{ctexart}
%       \input{assets/theme-violet}
%       \begin{document} ... \end{document}
% 只用 ctexart 自带能力 + xcolor/enumitem/fancyhdr/pifont/amsmath，
% 不依赖 tikz/tcolorbox，保证 MiKTeX 可直接编译。
%=====================================================================

\usepackage[a4paper,left=2.0cm,right=2.0cm,top=2.1cm,bottom=2.3cm,%
            headsep=12pt,footskip=22pt]{geometry}
\usepackage{amsmath,amssymb}
\usepackage{needspace}
\usepackage{xcolor}
\usepackage{array}
\usepackage{enumitem}
\usepackage{pifont}
\usepackage{fancyhdr}
\usepackage[hidelinks]{hyperref}

%---------------------------------------------------------------- 配色
\definecolor{violet}{RGB}{107,78,143}      % 主色：紫
\definecolor{violetdeep}{RGB}{63,44,92}    % 深紫：标题/强调
\definecolor{petal}{RGB}{183,159,208}      % 花瓣：装饰线、花饰
\definecolor{bandviolet}{RGB}{247,243,251} % 浅紫底：档案/提示框
\definecolor{bandcream}{RGB}{252,249,242}  % 米白底：易错/注意框
\definecolor{rulepurple}{RGB}{217,207,230} % 分隔线
\definecolor{inkgray}{RGB}{92,96,110}      % 次要文字

%---------------------------------------------------------------- 字体
% 有 Noto CJK 用 Noto，没有就退回 ctex 默认（Windows 下为宋体/黑体）
\IfFontExistsTF{Noto Serif CJK SC}{\setCJKmainfont{Noto Serif CJK SC}}{}
\IfFontExistsTF{Noto Sans CJK SC}{\setCJKsansfont{Noto Sans CJK SC}}{}

%---------------------------------------------------------------- 版式
\setlist{nosep,leftmargin=1.9em}
\linespread{1.28}
\setlength{\parskip}{3.5pt}
\setlength{\parindent}{0pt}
\setlength{\tabcolsep}{5pt}
\renewcommand{\arraystretch}{1.28}

%---------------------------------------------------------------- 花饰
% \ding{95} = 五瓣小花，\ding{96} = 星形花；都来自 pifont，必定存在
\newcommand{\orn}{\textcolor{petal}{\ding{95}}}          % 行内小标
\newcommand{\ornstar}{\textcolor{petal}{\ding{96}}}      % 星形花
\newcommand{\flowerrule}{%                                % 居中花饰分隔线
  \par\vspace{7pt}%
  \noindent\makebox[\linewidth]{\color{petal}%
    \rule{0.20\linewidth}{0.6pt}\hspace{8pt}{\small\ding{95}}\hspace{4pt}%
    \raisebox{1.1pt}{\tiny\ding{95}}\hspace{4pt}{\small\ding{95}}\hspace{8pt}%
    \rule{0.20\linewidth}{0.6pt}}%
  \par\vspace{7pt}}

%---------------------------------------------------------------- 标题与节头
\newcommand{\doctitle}[3]{%   #1 主标题  #2 副标题  #3 元信息（可用 \\ 分行）
  \begin{center}
    {\color{petal}\small \orn\enspace #2\enspace \ornstar}\\[7pt]
    {\LARGE\bfseries\color{violetdeep} #1}\\[10pt]
    {\color{petal}\rule{0.55\linewidth}{1pt}}\\[4pt]
    {\small\color{inkgray} #3}
  \end{center}\vspace{2pt}}

\newcommand{\hdr}[2]{%   #1 章节号（如 §0）  #2 标题
  \par\vspace{12pt}%
  \noindent\colorbox{violetdeep}{\color{white}\bfseries\small\;#1\;}\hspace{7pt}%
  {\color{violetdeep}\large\bfseries #2}\par\nobreak
  \vspace{3pt}{\color{petal}\hrule height 1.2pt}\vspace{6pt}}

\newcommand{\subhdr}[1]{%
  \par\vspace{6pt}%
  \noindent{\color{violet}\bfseries #1}\par\vspace{1pt}}

%---------------------------------------------------------------- 提示框
\newcommand{\callout}[1]{% 浅紫底：讲义档案、读法说明
  \par\medskip\noindent
  \fcolorbox{rulepurple}{bandviolet}{\parbox{\dimexpr\linewidth-2\fboxsep-2\fboxrule}{\vspace{2pt}#1\vspace{2pt}}}%
  \par\medskip}

\newcommand{\warnbox}[1]{% 米白底：易错、边界条件
  \par\medskip\noindent
  \fcolorbox{rulepurple}{bandcream}{\parbox{\dimexpr\linewidth-2\fboxsep-2\fboxrule}{\vspace{2pt}#1\vspace{2pt}}}%
  \par\medskip}

\newcommand{\gapitem}[2]{% #1 缺口类型（紫） #2 内容
  \par\vspace{2pt}\noindent
  \hangindent=1.4em\hangafter=1
  {\color{petal}\small\ding{96}}\enspace{\bfseries\color{violetdeep}#1}\enspace #2\par}

% 题干块（\qstem）：每道例题／作业的「详解」开头必须先复述题干与出处。
% 起因：2026-10-25 第五章讲义把七道作业的题面集中放在一页、详解另起一页，
% 读者翻到「第 4 题」时看不到条件，等于一道没有题干的题目。
% 规则：详解块内自带题干，读者不翻页就能读懂题在问什么。
\newcommand{\qstem}[1]{%
  \par\medskip\noindent
  \fcolorbox{violet}{bandviolet}{%
    \parbox{\dimexpr\linewidth-2\fboxsep-2\fboxrule}{%
      \vspace{2pt}%
      {\small\bfseries\color{violet}题干}\hspace{0.8em}#1%
      \vspace{2pt}}}%
  \par\medskip}

% v2: 保护题干与开头，长题可跨页，续页补上下文。
\newcommand{\qbreak}{\needspace{6\baselineskip}}
% 独立练习：ID、来源、题干、考点、答案标签（通常 ans:ID）。
\newcommand{\practice}[5]{%
  \qbreak\subhdr{练习 #1}\qstem{（#2）#3}%
  \textbf{考点}：#4\par
  详解见附录第~\pageref{#5}~页。\par\medskip}
% 附录开头：ID、来源、条件摘要；自动产生 ans:ID 标签。
\newcommand{\answerstart}[3]{%
  \qbreak\subhdr{#1\quad 详解}\label{ans:#1}\qstem{（#2）#3}}
\newcommand{\answercontinue}[2]{%
  \needspace{5\baselineskip}\subhdr{#1\quad 续解}%
  {\small\color{inkgray}本页接续：#2}\par\medskip}
% 延迟提供默认值，让旧讲义可在导言区保留自己的 \newcommand{\step}。
\AtBeginDocument{\providecommand{\step}[1]{\needspace{3\baselineskip}\par\medskip\textbf{#1}\par}}

% 前提旁注：只放在「需要理解的那一步」后面，给陈述 + 本题代入，不重证。
\newcommand{\aside}[1]{%
  \par\vspace{1pt}\noindent
  \hangindent=1.6em\hangafter=1
  {\small\color{inkgray}%
    \hspace*{0.6em}{\color{petal}$\triangleright$}\hspace{2pt}#1}%
  \par\vspace{2pt}}

%---------------------------------------------------------------- 表格
% 本机 MiKTeX 的 colortbl 与 array 版本不匹配（\rowcolor 会报错），
% 故表格统一用「上粗紫线 + 表头 + 中细线」的三线式，不依赖 colortbl。
\newcommand{\tabtop}{\par\vspace{3pt}{\color{violet}\hrule height 1.2pt}\vspace{1pt}}
\newcommand{\tabmid}{\par\vspace{1pt}{\color{petal}\hrule height 0.5pt}\vspace{1pt}}
\newcommand{\tabbot}{\par\vspace{3pt}{\color{rulepurple}\hrule height 0.5pt}\vspace{8pt}}
\newcommand{\thead}[1]{{\bfseries\color{violetdeep}#1}}

%---------------------------------------------------------------- 页眉页脚
\pagestyle{fancy}
\fancyhf{}
\renewcommand{\headrulewidth}{0.4pt}
\renewcommand{\headrule}{\hbox to\headwidth{\color{rulepurple}\leaders\hrule height \headrulewidth\hfill}}
\fancyhead[L]{\small\color{inkgray}\songti \doctag}
\fancyhead[R]{\small\color{inkgray}\songti \coursetag}
\fancyfoot[C]{\small\color{inkgray}\orn\enspace{\color{violet}\thepage}\enspace\orn}

% 每篇讲义在导言区设这两个值
\newcommand{\doctag}{}
\newcommand{\coursetag}{}
\newcommand{\settags}[2]{\renewcommand{\doctag}{#1}\renewcommand{\coursetag}{#2}}

%---------------------------------------------------------------- 数学记号
\newcommand{\vect}[1]{\boldsymbol{#1}}
\newcommand{\dd}{\,\mathrm{d}}
\newcommand{\ddx}{\frac{\mathrm{d}}{\mathrm{d}x}}
\addtolength{\jot}{2pt}



\setmainfont{Times New Roman}
\setsansfont{Times New Roman}
\usepackage{longtable}
\setlength{\headheight}{22pt}
\settags{俄语 · 名词总表}{俄语（第二外语）}
\newcommand{\zn}[1]{{\color{violetdeep}\bfseries #1}}
\begin{document}

\doctitle{俄语名词总表}
{按课件分类 · УРОК 1--8 · 中俄双语}
{每一课的名词都按课件的分法归类：先看性（он／она́／оно́），再看数（单数→复数）；\\
收录 \textbf{544} 条（含跨课复现），重音、词性、复数一并给出 \quad 2026-11-13}

\hdr{§0}{怎么用这份表}

\callout{\textbf{这份表解决什么问题}：期末考的不是「背词典」，是\textbf{拿到一个名词，能立刻说出它配 этот 还是 эта、配 мой 还是 моя}。所以全表只有一个组织原则：\textbf{按课件的分类走}——У4 按性分，У5 按复数词尾分，У6 按物主代词的四栏分，У7、У8 按形容词、指示代词与名词的搭配分。每一节标题后面都标了课件页码，「У4 p13」就是第 4 课 PDF 的第 13 页，可以随时回去对。}

\subhdr{0.1　三栏怎么读}

\begin{minipage}{\linewidth}
\tabtop
\begin{tabular}{@{}p{0.16\linewidth}p{0.30\linewidth}p{0.46\linewidth}@{}}
\thead{栏} & \thead{写法} & \thead{意思} \\
\end{tabular}
\tabmid
\begin{tabular}{@{}p{0.16\linewidth}p{0.30\linewidth}p{0.46\linewidth}@{}}
俄语 & сло́во、тетра́дь & 重音符「´」标在重读元音上；一个词通常只有一个重音 \\
中文 & 词、练习本 & 取课件语境里的常用义，不堆词典义项 \\
性 · 数 & 中 · слова́ & 「阳／阴／中」是性，「·」后面是复数形式；「复」表示只有复数 \\
\end{tabular}
\tabbot
\end{minipage}

\subhdr{0.2　四种特殊写法}

\begin{itemize}
\item \textbf{常用单数}：молоко́、му́зыка、пого́да 这类词平时不用复数，考试也不会让你变复数。
\item \textbf{不变格}：ко́фе、метро́、пальто́、такси́ 永远一个写法，性靠记（ко́фе 是阳性，其余多是中性）。
\item \textbf{单数 фрукт}：表里给的是课件上的复数形式 фру́кты，单数另注。
\item \textbf{\ding{72}}：出现在标题里，表示这一批是课件明确要求「需要记住／Запомните」的清单，优先级最高。
\end{itemize}

\warnbox{\textbf{一句提醒}：这份表是「认词」用的，不是「做题」用的。看的时候遮住右边两栏，只留俄语，先自己说中文和性；说不出来的，回到对应课的那一节，别在整张表上从头滑到尾。}

\subhdr{0.3　建议的三遍读法}

\begin{enumerate}
\item \textbf{第一遍（按课）}：从 §1 读到 §6，一课一课过。每节先看开头那句「这一课在讲什么」，再看表。
\item \textbf{第二遍（遮中文）}：只看俄语列，说中文、说性、说复数。错的词在页边画一道。
\item \textbf{第三遍（考前）}：只看 §7 速查卡和自己画过道的词；§7 是全表的「词尾规则压缩版」，考前一天翻两页就够。
\end{enumerate}

\begin{minipage}{\linewidth}
\subhdr{0.4　目录}
\tabtop
\begin{tabular}{@{}p{0.09\linewidth}p{0.76\linewidth}p{0.09\linewidth}@{}}
§1 & УРОК 1–3：语音课里先出现的名词 & \pageref{sec:1} \\
§2 & УРОК 4：名词的性（он／она́／оно́） & \pageref{sec:2} \\
§3 & УРОК 5：复数——三条规则与四张例外清单 & \pageref{sec:3} \\
§4 & УРОК 6：物主代词——按 мой／моя／моё／мои 四栏分 & \pageref{sec:4} \\
§5 & УРОК 7：形容词与名词——性、数必须一致 & \pageref{sec:5} \\
§6 & УРОК 8：指示代词与 нра́виться & \pageref{sec:6} \\
§7 & 考前速查：一张纸判性、变数 & \pageref{sec:7} \\
§8 & 补充：У8 之后课件里的新名词 & \pageref{sec:8} \\
\end{tabular}
\tabbot
\end{minipage}
\hdr{§1}{УРОК 1–3：语音课里先出现的名词}\label{sec:1}
У1–У3 是语音课，名词不是主角，只是练字母、重音、软硬辅音的材料。但这一批词是全课程最早出现的具体名词：先把它们认下来，У4 讲性、У5 讲复数时就不会被生词绊住。这三课的名词没有语法分类，就按课件自己的三张词表走。
\begin{longtable}{@{}p{0.27\linewidth}p{0.30\linewidth}p{0.37\linewidth}@{}}
\multicolumn{3}{@{}p{\linewidth}@{}}{\subhdr{1.1　У1 重音练习里的名词（课件 p28–p30）}\par{\small 课件用这些词练「重音落在哪个音节」。都是最具体的日常名词，先记形状和意思，性可以先不背——У4 会把它们重新分一次类。}} \\
\hline
\thead{俄语} & \thead{中文} & \thead{性 · 复数} \\
\hline
\endfirsthead
\hline
\thead{俄语} & \thead{中文} & \thead{性 · 复数} \\
\multicolumn{3}{@{}l}{\scriptsize\color{inkgray}（接上页）1.1　У1 重音练习里的名词（课件 p28–p30）} \\
\hline
\endhead
дом & 房子 & 阳 · дома́ \\
кот & 猫 & 阳 · коты́ \\
стол & 桌子 & 阳 · столы́ \\
сыр & 奶酪 & 阳 · сы́ры \\
ка́рта & 地图 & 阴 · ка́рты \\
бу́ква & 字母 & 阴 · бу́квы \\
па́нда & 熊猫 & 阴 · па́нды \\
ма́ма & 妈妈 & 阴 · ма́мы \\
со́ус & 酱汁 & 阳 · со́усы \\
ры́ба & 鱼 & 阴 · ры́бы \\
фру́кты & 水果 & 复 · 单数 фрукт \\
бана́н & 香蕉 & 阳 · бана́ны \\
уро́к & 课 & 阳 · уро́ки \\
му́зыка & 音乐 & 阴 · 常用单数 \\
суббо́та & 星期六 & 阴 · суббо́ты \\
вода́ & 水 & 阴 · во́ды \\
нога́ & 腿、脚 & 阴 · но́ги \\
молоко́ & 牛奶 & 中 · 常用单数 \\
фо́то & 照片 & 中 · 不变格 \\
\hline
\end{longtable}
\begin{longtable}{@{}p{0.27\linewidth}p{0.30\linewidth}p{0.37\linewidth}@{}}
\multicolumn{3}{@{}p{\linewidth}@{}}{\subhdr{1.2　У2 软硬辅音词表（课件 p9–p10、p30）}\par{\small 这一课练的是「辅音在软音符号、е／ё／ю／я／и 前变软」。词表本身按语音分组，这里仍按词义排；带 -ь 的词（день、мать、вещь…）正是 У4「-ь 判不出性」那两栏的提前预演。}} \\
\hline
\thead{俄语} & \thead{中文} & \thead{性 · 复数} \\
\hline
\endfirsthead
\hline
\thead{俄语} & \thead{中文} & \thead{性 · 复数} \\
\multicolumn{3}{@{}l}{\scriptsize\color{inkgray}（接上页）1.2　У2 软硬辅音词表（课件 p9–p10、p30）} \\
\hline
\endhead
мел & 粉笔 & 阳 · мелки́ \\
ме́сто & 地方 & 中 · места́ \\
вещь & 东西 & 阴 · ве́щи \\
текст & 课文 & 阳 · те́ксты \\
день & 白天、日子 & 阳 · дни \\
метро́ & 地铁 & 中 · 不变格 \\
меню́ & 菜单 & 中 · 不变格 \\
сестра́ & 姐妹 & 阴 · сёстры \\
семья́ & 家庭 & 阴 · се́мьи \\
весна́ & 春天 & 阴 · вёсны \\
мать & 母亲 & 阴 · ма́тери \\
лес & 森林 & 阳 · леса́ \\
тётя & 阿姨、姑姑 & 阴 · тёти \\
дя́дя & 叔叔、舅舅 & 阳 · дя́ди \\
пе́сня & 歌 & 阴 · пе́сни \\
река́ & 河 & 阴 · ре́ки \\
оте́ц & 父亲 & 阳 · отцы́ \\
друг & 朋友 & 阳 · друзья́ \\
студе́нт & 大学生 & 阳 · студе́нты \\
студе́нтка & 女大学生 & 阴 · студе́нтки \\
апте́ка & 药店 & 阴 · апте́ки \\
карти́на & 画 & 阴 · карти́ны \\
туале́т & 厕所 & 阳 · туале́ты \\
лифт & 电梯 & 阳 · ли́фты \\
кни́га & 书 & 阴 · кни́ги \\
аудито́рия & 教室 & 阴 · аудито́рии \\
преподава́тель & 教师 & 阳 · преподава́тели \\
университе́т & 大学 & 阳 · университе́ты \\
улы́бка & 微笑 & 阴 · улы́бки \\
вто́рник & 星期二 & 阳 · вто́рники \\
остано́вка & 车站 & 阴 · остано́вки \\
вход & 入口 & 阳 · вхо́ды \\
ло́жка & 勺子 & 阴 · ло́жки \\
\hline
\end{longtable}
\begin{longtable}{@{}p{0.27\linewidth}p{0.30\linewidth}p{0.37\linewidth}@{}}
\multicolumn{3}{@{}p{\linewidth}@{}}{\subhdr{1.3　У3 食物与日常词表（课件 p7、p16–p21）}\par{\small У3 的语法是 да／нет 问答和 ИК-3 语调，词表集中在「吃什么」和「在哪儿」。食物类词大多是中性或阳性，阴性只有 шко́ла 一类 -а 结尾的。}} \\
\hline
\thead{俄语} & \thead{中文} & \thead{性 · 复数} \\
\hline
\endfirsthead
\hline
\thead{俄语} & \thead{中文} & \thead{性 · 复数} \\
\multicolumn{3}{@{}l}{\scriptsize\color{inkgray}（接上页）1.3　У3 食物与日常词表（课件 p7、p16–p21）} \\
\hline
\endhead
мя́со & 肉 & 中 · 常用单数 \\
сок & 果汁 & 阳 · со́ки \\
хлеб & 面包 & 阳 · хлеба́ \\
рис & 大米 & 阳 · 常用单数 \\
чи́псы & 薯片 & 复 · чи́псы \\
то́фу & 豆腐 & 中 · 不变格 \\
чай & 茶 & 阳 · 常用单数 \\
па́па & 爸爸 & 阳 · па́пы \\
ру́чка & 笔 & 阴 · ру́чки \\
каранда́ш & 铅笔 & 阳 · карандаши́ \\
магази́н & 商店 & 阳 · магази́ны \\
банк & 银行 & 阳 · ба́нки \\
за́втрак & 早饭 & 阳 · за́втраки \\
обе́д & 午饭 & 阳 · обе́ды \\
у́жин & 晚饭 & 阳 · у́жины \\
\hline
\end{longtable}
\hdr{§2}{УРОК 4：名词的性（он／она́／оно́）}\label{sec:2}
У4 只做一件事：给名词定一个性。课件把它压成一张表（p11–p12）——\textbf{阳性}看硬辅音、-й、-ь，\textbf{阴性}看 -а、-я、-ь，\textbf{中性}看 -о、-е、-мя；然后 p13 单独把 -ь 结尾的名词分成两栏，意思很直白：\textbf{-ь 判不出来，只能整词记}。下面先把这张表抄清，再按三性把 У4 出现过的词全部排进去。
\needspace{6\baselineskip}\subhdr{2.1　课件原表：三个性与词尾（p11–p12）}

\begin{minipage}{\linewidth}
\tabtop
\begin{tabular}{@{}p{0.08\linewidth}p{0.08\linewidth}p{0.22\linewidth}p{0.56\linewidth}@{}}
\thead{性} & \thead{代词} & \thead{词尾} & \thead{课件例词} \\
\end{tabular}
\tabmid
\begin{tabular}{@{}p{0.08\linewidth}p{0.08\linewidth}p{0.22\linewidth}p{0.56\linewidth}@{}}
阳性 & он & 硬辅音、-й、-ь & дом, музе́й, слова́рь \\
阴性 & она́ & -а、-я、-ь & ко́мната, пе́сня, тетра́дь \\
中性 & оно́ & -о、-е、-мя & окно́, мо́ре, вре́мя \\
例外 & он & -а／-я & па́па, де́душка, дя́дя, мужчи́на, ю́ноша \\
\end{tabular}
\tabbot
\end{minipage}

\textbf{这张表只有一处「脏」}：-ь 同时挂在阳性和阴性两行，光看词尾判不出来。所以 p13 给了两栏各 19 个词（下面带 \ding{72} 的两节），意思是「这一格没有规律，只能整词记」。判断顺序永远是：\textbf{先看词尾能不能直接判}（-а／-я → 阴，-о／-е → 中，辅音 → 阳），\textbf{判不了才查 -ь 两栏}。
\begin{longtable}{@{}p{0.27\linewidth}p{0.30\linewidth}p{0.37\linewidth}@{}}
\multicolumn{3}{@{}p{\linewidth}@{}}{\subhdr{2.2　阳性：硬辅音结尾（课件 p10、p14–p16）}\par{\small 这一批是阳性里最「干净」的：词尾是 б、в、г、д、ж、з、к、л、м、н、п、р、с、т、ф、х、ц、ч、ш、щ 之一，直接配 он。}} \\
\hline
\thead{俄语} & \thead{中文} & \thead{性 · 复数} \\
\hline
\endfirsthead
\hline
\thead{俄语} & \thead{中文} & \thead{性 · 复数} \\
\multicolumn{3}{@{}l}{\scriptsize\color{inkgray}（接上页）2.2　阳性：硬辅音结尾（课件 p10、p14–p16）} \\
\hline
\endhead
дом & 房子 & 阳 · дома́ \\
парк & 公园 & 阳 · па́рки \\
заво́д & 工厂 & 阳 · заво́ды \\
брат & 兄弟 & 阳 · бра́тья \\
стол & 桌子 & 阳 · столы́ \\
стул & 椅子 & 阳 · сту́лья \\
друг & 朋友 & 阳 · друзья́ \\
суп & 汤 & 阳 · су́пы \\
го́род & 城市 & 阳 · города́ \\
сын & 儿子 & 阳 · сыновья́ \\
кот & 猫 & 阳 · коты́ \\
банк & 银行 & 阳 · ба́нки \\
авто́бус & 公共汽车 & 阳 · авто́бусы \\
расска́з & 故事 & 阳 · расска́зы \\
сыр & 奶酪 & 阳 · сы́ры \\
торт & 蛋糕 & 阳 · то́рты \\
студе́нт & 大学生 & 阳 · студе́нты \\
зонт & 伞 & 阳 · зо́нты \\
лес & 森林 & 阳 · леса́ \\
врач & 医生 & 阳 · врачи́ \\
океа́н & 海洋 & 阳 · океа́ны \\
класс & 班级 & 阳 · кла́ссы \\
час & 小时 & 阳 · часы́ \\
звук & 声音 & 阳 · зву́ки \\
слог & 音节 & 阳 · сло́ги \\
сад & 花园 & 阳 · сады́ \\
гара́ж & 车库 & 阳 · гаражи́ \\
альбо́м & 相册 & 阳 · альбо́мы \\
телефо́н & 电话 & 阳 · телефо́ны \\
биле́т & 票 & 阳 · биле́ты \\
муж & 丈夫 & 阳 · мужья́ \\
сувени́р & 纪念品 & 阳 · сувени́ры \\
сок & 果汁 & 阳 · со́ки \\
фило́лог & 语文学家 & 阳 · фило́логи \\
бизнесме́н & 商人 & 阳 · бизнесме́ны \\
диало́г & 对话 & 阳 · диало́ги \\
арти́ст & 演员 & 阳 · арти́сты \\
био́лог & 生物学家 & 阳 · био́логи \\
конве́рт & 信封 & 阳 · конве́рты \\
глаз & 眼睛 & 阳 · глаза́ \\
журнали́ст & 记者 & 阳 · журнали́сты \\
экономи́ст & 经济学家 & 阳 · экономи́сты \\
до́ктор & 医生 & 阳 · доктора́ \\
това́рищ & 同志、同学 & 阳 · това́рищи \\
обе́д & 午饭 & 阳 · обе́ды \\
за́втрак & 早饭 & 阳 · за́втраки \\
у́жин & 晚饭 & 阳 · у́жины \\
а́дрес & 地址 & 阳 · адреса́ \\
го́лос & 嗓音 & 阳 · голоса́ \\
инжене́р & 工程师 & 阳 · инжене́ры \\
ве́чер & 晚上 & 阳 · вечера́ \\
пра́здник & 节日 & 阳 · пра́здники \\
центр & 中心 & 阳 · це́нтры \\
\hline
\end{longtable}
\begin{longtable}{@{}p{0.27\linewidth}p{0.30\linewidth}p{0.37\linewidth}@{}}
\multicolumn{3}{@{}p{\linewidth}@{}}{\subhdr{2.3　阳性：-й 结尾（课件 p12、p15）}\par{\small -й 结尾一律阳性，复数把 -й 换成 -и：музе́й → музе́и。同一组里还有 санато́рий（У5 p6），复数 санато́рии。}} \\
\hline
\thead{俄语} & \thead{中文} & \thead{性 · 复数} \\
\hline
\endfirsthead
\hline
\thead{俄语} & \thead{中文} & \thead{性 · 复数} \\
\multicolumn{3}{@{}l}{\scriptsize\color{inkgray}（接上页）2.3　阳性：-й 结尾（课件 p12、p15）} \\
\hline
\endhead
музе́й & 博物馆 & 阳 · музе́и \\
геро́й & 英雄 & 阳 · геро́и \\
санато́рий & 疗养院 & 阳 · санато́рии \\
\hline
\end{longtable}
\begin{longtable}{@{}p{0.27\linewidth}p{0.30\linewidth}p{0.37\linewidth}@{}}
\multicolumn{3}{@{}p{\linewidth}@{}}{\subhdr{\ding{72} 2.4　阳性：-ь 结尾（课件 p13「需要记住」19 个 ＋ p16）}\par{\small 这 19 个是课件 p13 左栏原表，顺序照抄；最后一个 гость、писа́тель 来自 p16 的分栏练习。它们和下一节阴性那 19 个**必须成对背**：同样的 -ь，一半阳一半阴，没有任何词尾线索。}} \\
\hline
\thead{俄语} & \thead{中文} & \thead{性 · 复数} \\
\hline
\endfirsthead
\hline
\thead{俄语} & \thead{中文} & \thead{性 · 复数} \\
\multicolumn{3}{@{}l}{\scriptsize\color{inkgray}（接上页）\ding{72} 2.4　阳性：-ь 结尾（课件 p13「需要记住」19 个 ＋ p16）} \\
\hline
\endhead
слова́рь & 词典 & 阳 · словари́ \\
преподава́тель & 教师 & 阳 · преподава́тели \\
день & 白天、日子 & 阳 · дни \\
календа́рь & 日历 & 阳 · календари́ \\
портфе́ль & 公文包 & 阳 · портфе́ли \\
карто́фель & 土豆 & 阳 · 常用单数 \\
автомоби́ль & 汽车 & 阳 · автомоби́ли \\
дождь & 雨 & 阳 · дожди́ \\
ка́шель & 咳嗽 & 阳 · ка́шли \\
кора́бль & 船 & 阳 · корабли́ \\
ли́вень & 暴雨 & 阳 · ли́вни \\
ло́коть & 肘 & 阳 · ло́кти \\
но́готь & 指甲 & 阳 · но́гти \\
ноль & 零 & 阳 · нули́ \\
ого́нь & 火 & 阳 · огни́ \\
по́лдень & 正午 & 阳 · по́лдни \\
реме́нь & 皮带 & 阳 · ремни́ \\
юа́нь & 元（人民币） & 阳 · юа́ни \\
спекта́кль & 戏剧 & 阳 · спекта́кли \\
гость & 客人 & 阳 · го́сти \\
писа́тель & 作家 & 阳 · писа́тели \\
\hline
\end{longtable}
\begin{longtable}{@{}p{0.27\linewidth}p{0.30\linewidth}p{0.37\linewidth}@{}}
\multicolumn{3}{@{}p{\linewidth}@{}}{\subhdr{2.5　阴性：-а／-я 结尾（课件 p10、p12、p14–p16）}\par{\small 阴性主力。注意 пого́да、му́зыка 这类「只有单数」的词也在这里，性照样是阴性，只是不用复数。}} \\
\hline
\thead{俄语} & \thead{中文} & \thead{性 · 复数} \\
\hline
\endfirsthead
\hline
\thead{俄语} & \thead{中文} & \thead{性 · 复数} \\
\multicolumn{3}{@{}l}{\scriptsize\color{inkgray}（接上页）2.5　阴性：-а／-я 结尾（课件 p10、p12、p14–p16）} \\
\hline
\endhead
ко́мната & 房间 & 阴 · ко́мнаты \\
ла́мпа & 灯 & 阴 · ла́мпы \\
пе́сня & 歌 & 阴 · пе́сни \\
аудито́рия & 教室 & 阴 · аудито́рии \\
подру́га & 女性朋友 & 阴 · подру́ги \\
гру́ппа & 小组、班级 & 阴 · гру́ппы \\
ры́ба & 鱼 & 阴 · ры́бы \\
бу́ква & 字母 & 阴 · бу́квы \\
страна́ & 国家 & 阴 · стра́ны \\
соба́ка & 狗 & 阴 · соба́ки \\
су́мка & 包 & 阴 · су́мки \\
вода́ & 水 & 阴 · во́ды \\
сестра́ & 姐妹 & 阴 · сёстры \\
ма́ма & 妈妈 & 阴 · ма́мы \\
ба́бушка & 奶奶、外婆 & 阴 · ба́бушки \\
жена́ & 妻子 & 阴 · жёны \\
студе́нтка & 女大学生 & 阴 · студе́нтки \\
дере́вня & 村庄 & 阴 · дере́вни \\
река́ & 河 & 阴 · ре́ки \\
мину́та & 分钟 & 阴 · мину́ты \\
ша́пка & 帽子 & 阴 · ша́пки \\
доро́га & 路 & 阴 · доро́ги \\
экску́рсия & 游览 & 阴 · экску́рсии \\
ва́за & 花瓶 & 阴 · ва́зы \\
пого́да & 天气 & 阴 · 常用单数 \\
кварти́ра & 住宅 & 阴 · кварти́ры \\
ци́фра & 数字 & 阴 · ци́фры \\
ста́нция & 车站 & 阴 · ста́нции \\
столи́ца & 首都 & 阴 · столи́цы \\
газе́та & 报纸 & 阴 · газе́ты \\
ма́рка & 邮票 & 阴 · ма́рки \\
домохозя́йка & 家庭主妇 & 阴 · домохозя́йки \\
страни́ца & 页 & 阴 · страни́цы \\
ку́ртка & 夹克 & 阴 · ку́ртки \\
шу́ба & 皮大衣 & 阴 · шу́бы \\
\hline
\end{longtable}
\begin{longtable}{@{}p{0.27\linewidth}p{0.30\linewidth}p{0.37\linewidth}@{}}
\multicolumn{3}{@{}p{\linewidth}@{}}{\subhdr{\ding{72} 2.6　阴性：-ь 结尾（课件 p13「需要记住」19 个 ＋ p15）}\par{\small 与 2.4 的 19 个一一对应（另外多一个 вещь，来自 p15 的分栏练习）。最容易混的四对：день（阳）／ночь（阴）、дождь（阳）／тетра́дь（阴）、портфе́ль（阳）／посте́ль（阴）、юа́нь（阳）／меда́ль（阴）。背的时候成对背。}} \\
\hline
\thead{俄语} & \thead{中文} & \thead{性 · 复数} \\
\hline
\endfirsthead
\hline
\thead{俄语} & \thead{中文} & \thead{性 · 复数} \\
\multicolumn{3}{@{}l}{\scriptsize\color{inkgray}（接上页）\ding{72} 2.6　阴性：-ь 结尾（课件 p13「需要记住」19 个 ＋ p15）} \\
\hline
\endhead
тетра́дь & 练习本 & 阴 · тетра́ди \\
пло́щадь & 广场 & 阴 · пло́щади \\
ночь & 夜 & 阴 · но́чи \\
дверь & 门 & 阴 · две́ри \\
о́сень & 秋天 & 阴 · о́сени \\
крова́ть & 床 & 阴 · крова́ти \\
но́вость & 新闻（一条） & 阴 · но́вости \\
ме́бель & 家具 & 阴 · 常用单数 \\
о́бувь & 鞋类 & 阴 · 常用单数 \\
любо́вь & 爱 & 阴 · 常用单数 \\
ладо́нь & 手掌 & 阴 · ладо́ни \\
мете́ль & 暴风雪 & 阴 · мете́ли \\
мысль & 思想 & 阴 · мы́сли \\
о́чередь & 队、轮次 & 阴 · очереди \\
па́мять & 记忆 & 阴 · 常用单数 \\
посте́ль & 床铺 & 阴 · посте́ли \\
цель & 目标 & 阴 · це́ли \\
по́лночь & 午夜 & 阴 · 常用单数 \\
меда́ль & 奖章 & 阴 · меда́ли \\
вещь & 东西 & 阴 · ве́щи \\
\hline
\end{longtable}
\begin{longtable}{@{}p{0.27\linewidth}p{0.30\linewidth}p{0.37\linewidth}@{}}
\multicolumn{3}{@{}p{\linewidth}@{}}{\subhdr{2.7　中性：-о／-е 结尾（课件 p12、p14–p16）}\par{\small 中性最省事：-о、-е 结尾直接配 оно́。注意 фо́то、пальто́ 是外来词，属于「不变格」，性靠「它」来判断。}} \\
\hline
\thead{俄语} & \thead{中文} & \thead{性 · 复数} \\
\hline
\endfirsthead
\hline
\thead{俄语} & \thead{中文} & \thead{性 · 复数} \\
\multicolumn{3}{@{}l}{\scriptsize\color{inkgray}（接上页）2.7　中性：-о／-е 结尾（课件 p12、p14–p16）} \\
\hline
\endhead
окно́ & 窗户 & 中 · о́кна \\
молоко́ & 牛奶 & 中 · 常用单数 \\
мо́ре & 海 & 中 · моря́ \\
зда́ние & 建筑 & 中 · зда́ния \\
фо́то & 照片 & 中 · 不变格 \\
со́лнце & 太阳 & 中 · 常用单数 \\
о́зеро & 湖 & 中 · озёра \\
упражне́ние & 练习 & 中 · упражне́ния \\
пальто́ & 大衣 & 中 · 不变格 \\
по́ле & 田野 & 中 · поля́ \\
письмо́ & 信 & 中 · пи́сьма \\
лицо́ & 脸 & 中 · ли́ца \\
ме́сто & 地方 & 中 · места́ \\
предложе́ние & 句子 & 中 · предложе́ния \\
сло́во & 词 & 中 · слова́ \\
ударе́ние & 重音 & 中 · ударе́ния \\
пла́тье & 连衣裙 & 中 · пла́тья \\
общежи́тие & 宿舍 & 中 · общежи́тия \\
\hline
\end{longtable}
\begin{longtable}{@{}p{0.27\linewidth}p{0.30\linewidth}p{0.37\linewidth}@{}}
\multicolumn{3}{@{}p{\linewidth}@{}}{\subhdr{2.8　中性：-мя 结尾（课件 p12、p16）}\par{\small 这一组只有几个词，但每个都要背：-мя 结尾全是中性，复数把 -мя 换成 -мена。}} \\
\hline
\thead{俄语} & \thead{中文} & \thead{性 · 复数} \\
\hline
\endfirsthead
\hline
\thead{俄语} & \thead{中文} & \thead{性 · 复数} \\
\multicolumn{3}{@{}l}{\scriptsize\color{inkgray}（接上页）2.8　中性：-мя 结尾（课件 p12、p16）} \\
\hline
\endhead
вре́мя & 时间 & 中 · времена́ \\
и́мя & 名字 & 中 · имена́ \\
\hline
\end{longtable}
\begin{longtable}{@{}p{0.27\linewidth}p{0.30\linewidth}p{0.37\linewidth}@{}}
\multicolumn{3}{@{}p{\linewidth}@{}}{\subhdr{2.9　例外：-а／-я 结尾却是阳性（课件 p10、p12）}\par{\small 这五个都是「指人的 -а／-я」，是阳性里唯一需要单独记的一格。数量很少，别把这条推广成「-а 也可能是阳性」——那样会写错绝大多数词。}} \\
\hline
\thead{俄语} & \thead{中文} & \thead{性 · 复数} \\
\hline
\endfirsthead
\hline
\thead{俄语} & \thead{中文} & \thead{性 · 复数} \\
\multicolumn{3}{@{}l}{\scriptsize\color{inkgray}（接上页）2.9　例外：-а／-я 结尾却是阳性（课件 p10、p12）} \\
\hline
\endhead
па́па & 爸爸 & 阳 · па́пы \\
де́душка & 爷爷、外公 & 阳 · де́душки \\
дя́дя & 叔叔、舅舅 & 阳 · дя́ди \\
мужчи́на & 男人 & 阳 · мужчи́ны \\
ю́ноша & 青年 & 阳 · ю́ноши \\
\hline
\end{longtable}
\hdr{§3}{УРОК 5：复数——三条规则与四张例外清单}\label{sec:3}
У5 把复数拆成三条规则（阳性、阴性、中性各一条），但真正要背的其实是后面的四张清单：\textbf{特殊复数}、\textbf{只有单数}、\textbf{只有复数}、\textbf{不变格}。最后五张主题词表（城市／大学／课堂／宿舍／家庭）是课件的写作素材，组句题多半从这里取词。
\needspace{6\baselineskip}\subhdr{3.1　三条规则（课件 p6–p8）}

\begin{minipage}{\linewidth}
\tabtop
\begin{tabular}{@{}p{0.06\linewidth}p{0.27\linewidth}p{0.08\linewidth}p{0.50\linewidth}@{}}
\thead{性} & \thead{单数词尾} & \thead{复数} & \thead{课件例词} \\
\end{tabular}
\tabmid
\begin{tabular}{@{}p{0.06\linewidth}p{0.27\linewidth}p{0.08\linewidth}p{0.50\linewidth}@{}}
阳 & 硬辅音 & -ы & студе́нт → студе́нты；университе́т → университе́ты \\
阳 & -г／-к／-х／-ж／\newline -ш／-ч／-щ & -и & слог → сло́ги；уро́к → уро́ки；нож → ножи́ \\
阳 & -ь／-й & -и & слова́рь → словари́；музе́й → музе́и \\
阴 & -а & -ы & сестра́ → сёстры；газе́та → газе́ты \\
阴 & -га／-ка／-ха／\newline -жа／-ша／-ча／-ща & -и & кни́га → кни́ги；студе́нтка → студе́нтки \\
阴 & -ь & -и & тетра́дь → тетра́ди \\
阴 & -я／-ия & -и & пе́сня → пе́сни；аудито́рия → аудито́рии \\
中 & -о & -а & окно́ → о́кна \\
中 & -е／-ие & -я／-ия & по́ле → поля́；упражне́ние → упражне́ния \\
\end{tabular}
\tabbot
\end{minipage}

\textbf{八个字}：\textbf{硬词干加 ы／а，软词干加 и／я}。第二行和第五行是同一条规则的两个性别版本：\textbf{г、к、х、ж、ш、ч、щ 后面永远写 и，不写 ы}。第三行和第六行也是同一条：\textbf{-ь 一律变 -и}。
\begin{longtable}{@{}p{0.22\linewidth}p{0.30\linewidth}p{0.44\linewidth}@{}}
\multicolumn{3}{@{}p{\linewidth}@{}}{\subhdr{\ding{72} 3.2　特殊复数：24 组不能按规则推的词（课件 p18）}\par{\small 这 24 组在 У7 p18 又完整出现一次，是课件唯一重复列出的词表，重要程度可见一斑。它们分三类：\textbf{词干变了}（мать → ма́тери、и́мя → имена́、ребёнок → де́ти）、\textbf{重音跑到词尾}（дом → дома́、го́род → города́）、\textbf{多出 -ья-}（брат → бра́тья、стул → сту́лья）。读的时候连着重音一起读。}} \\
\hline
\thead{单数} & \thead{复数} & \thead{中文} \\
\hline
\endfirsthead
\hline
\thead{单数} & \thead{复数} & \thead{中文} \\
\multicolumn{3}{@{}l}{\scriptsize\color{inkgray}（接上页）\ding{72} 3.2　特殊复数：24 组不能按规则推的词（课件 p18）} \\
\hline
\endhead
мать（阴） & ма́тери & 母亲 \\
дочь（阴） & до́чери & 女儿 \\
сестра́（阴） & сёстры & 姐妹 \\
оте́ц（阳） & отцы́ & 父亲 \\
день（阳） & дни & 白天、日子 \\
и́мя（中） & имена́ & 名字 \\
вре́мя（中） & времена́ & 时间 \\
я́блоко（中） & я́блоки & 苹果 \\
ребёнок（阳） & де́ти & 孩子 \\
челове́к（阳） & лю́ди & 人 \\
дом（阳） & дома́ & 房子 \\
го́род（阳） & города́ & 城市 \\
бе́рег（阳） & берега́ & 岸 \\
лес（阳） & леса́ & 森林 \\
о́стров（阳） & острова́ & 岛 \\
по́езд（阳） & поезда́ & 火车 \\
ве́чер（阳） & вечера́ & 晚上 \\
брат（阳） & бра́тья & 兄弟 \\
сын（阳） & сыновья́ & 儿子 \\
друг（阳） & друзья́ & 朋友 \\
муж（阳） & мужья́ & 丈夫 \\
де́рево（中） & дере́вья & 树 \\
лист（阳） & ли́стья & 叶子 \\
стул（阳） & сту́лья & 椅子 \\
\hline
\end{longtable}
\begin{longtable}{@{}p{0.27\linewidth}p{0.30\linewidth}p{0.37\linewidth}@{}}
\multicolumn{3}{@{}p{\linewidth}@{}}{\subhdr{3.3　只有单数（课件 p9）}\par{\small 这些词没有复数形式。考试里如果让你变复数，它们要么不考，要么答案只能是「这个词不用复数」。}} \\
\hline
\thead{俄语} & \thead{中文} & \thead{性 · 复数} \\
\hline
\endfirsthead
\hline
\thead{俄语} & \thead{中文} & \thead{性 · 复数} \\
\multicolumn{3}{@{}l}{\scriptsize\color{inkgray}（接上页）3.3　只有单数（课件 p9）} \\
\hline
\endhead
со́лнце & 太阳 & 中 · 常用单数 \\
луна́ & 月亮 & 阴 · 常用单数 \\
молоко́ & 牛奶 & 中 · 常用单数 \\
мя́со & 肉 & 中 · 常用单数 \\
рис & 大米 & 阳 · 常用单数 \\
карто́фель & 土豆 & 阳 · 常用单数 \\
лапша́ & 面条 & 阴 · 常用单数 \\
лук & 洋葱 & 阳 · 常用单数 \\
\hline
\end{longtable}
\begin{longtable}{@{}p{0.27\linewidth}p{0.30\linewidth}p{0.37\linewidth}@{}}
\multicolumn{3}{@{}p{\linewidth}@{}}{\subhdr{3.4　只有复数（课件 p10）}\par{\small 「只有复数」不等于「复数难写」，而是\textbf{它没有单数可用}：часы́ 指整件「钟表」，要说「一只表」不用单数。这类词配形容词、代词时一律用复数（э́ти часы́、мои́ де́ньги）。}} \\
\hline
\thead{俄语} & \thead{中文} & \thead{性 · 复数} \\
\hline
\endfirsthead
\hline
\thead{俄语} & \thead{中文} & \thead{性 · 复数} \\
\multicolumn{3}{@{}l}{\scriptsize\color{inkgray}（接上页）3.4　只有复数（课件 p10）} \\
\hline
\endhead
роди́тели & 父母 & 复 · роди́тели \\
де́ньги & 钱 & 复 · де́ньги \\
часы́ & 钟表 & 复 · часы́ \\
очки́ & 眼镜 & 复 · очки́ \\
ножни́цы & 剪刀 & 复 · ножни́цы \\
брю́ки & 裤子 & 复 · брю́ки \\
джи́нсы & 牛仔裤 & 复 · джи́нсы \\
духи́ & 香水 & 复 · духи́ \\
макаро́ны & 通心粉 & 复 · макаро́ны \\
консе́рвы & 罐头 & 复 · консе́рвы \\
сли́вки & 奶油 & 复 · сли́вки \\
щи & 白菜汤 & 复 · щи \\
канику́лы & 假期 & 复 · канику́лы \\
су́тки & 一昼夜 & 复 · су́тки \\
воро́та & 大门 & 复 · воро́та \\
весы́ & 秤 & 复 · весы́ \\
кавы́чки & 引号 & 复 · кавы́чки \\
\hline
\end{longtable}
\begin{longtable}{@{}p{0.27\linewidth}p{0.30\linewidth}p{0.37\linewidth}@{}}
\multicolumn{3}{@{}p{\linewidth}@{}}{\subhdr{3.5　不变格名词：永远一个写法（课件 p10）}\par{\small 外来词，性靠单独记：\textbf{ко́фе 是阳性}（он ко́фе），其余绝大多数是中性。考试里它们最容易在「配 этот 还是 эта」上出错——ко́фе 配 э́тот，кафе́ 配 э́то。}} \\
\hline
\thead{俄语} & \thead{中文} & \thead{性 · 复数} \\
\hline
\endfirsthead
\hline
\thead{俄语} & \thead{中文} & \thead{性 · 复数} \\
\multicolumn{3}{@{}l}{\scriptsize\color{inkgray}（接上页）3.5　不变格名词：永远一个写法（课件 p10）} \\
\hline
\endhead
ко́фе & 咖啡 & 阳 · 不变格 \\
кафе́ & 咖啡馆 & 中 · 不变格 \\
метро́ & 地铁 & 中 · 不变格 \\
пальто́ & 大衣 & 中 · 不变格 \\
ра́дио & 收音机 & 中 · 不变格 \\
пиани́но & 钢琴 & 中 · 不变格 \\
такси́ & 出租车 & 中 · 不变格 \\
шоссе́ & 公路 & 中 · 不变格 \\
кака́о & 可可 & 中 · 不变格 \\
интервью́ & 采访 & 中 · 不变格 \\
ателье́ & 裁缝店 & 中 · 不变格 \\
бюро́ & 局、办事处 & 中 · 不变格 \\
фойе́ & 休息厅 & 中 · 不变格 \\
\hline
\end{longtable}
\begin{longtable}{@{}p{0.27\linewidth}p{0.30\linewidth}p{0.37\linewidth}@{}}
\multicolumn{3}{@{}p{\linewidth}@{}}{\subhdr{3.6　主题词表：城市（课件 p11、p5）}\par{\small 课件 p11 的写作模型是 \textbf{Э́то го́род. Здесь проспе́кты, у́лицы, пло́щади…}（这是城市，这里有大街、街道、广场……），要求把每个名词都变成复数。所以这张表要\textbf{横着背}：单数—复数一起记。}} \\
\hline
\thead{俄语} & \thead{中文} & \thead{性 · 复数} \\
\hline
\endfirsthead
\hline
\thead{俄语} & \thead{中文} & \thead{性 · 复数} \\
\multicolumn{3}{@{}l}{\scriptsize\color{inkgray}（接上页）3.6　主题词表：城市（课件 p11、p5）} \\
\hline
\endhead
проспе́кт & 大街 & 阳 · проспе́кты \\
у́лица & 街道 & 阴 · у́лицы \\
пло́щадь & 广场 & 阴 · пло́щади \\
университе́т & 大学 & 阳 · университе́ты \\
поликли́ника & 诊所 & 阴 · поликли́ники \\
кафе́ & 咖啡馆 & 中 · 不变格 \\
апте́ка & 药店 & 阴 · апте́ки \\
общежи́тие & 宿舍 & 中 · общежи́тия \\
магази́н & 商店 & 阳 · магази́ны \\
музе́й & 博物馆 & 阳 · музе́и \\
теа́тр & 剧院 & 阳 · теа́тры \\
вы́ставка & 展览 & 阴 · вы́ставки \\
суперма́ркет & 超市 & 阳 · суперма́ркеты \\
шко́ла & 学校 & 阴 · шко́лы \\
посо́льство & 大使馆 & 中 · посо́льства \\
вокза́л & 火车站 & 阳 · вокза́лы \\
банк & 银行 & 阳 · ба́нки \\
заво́д & 工厂 & 阳 · заво́ды \\
фа́брика & 工厂 & 阴 · фа́брики \\
кио́ск & 报刊亭 & 阳 · кио́ски \\
кинотеа́тр & 电影院 & 阳 · кинотеа́тры \\
рестора́н & 餐厅 & 阳 · рестора́ны \\
гости́ница & 旅馆 & 阴 · гости́ницы \\
дом & 房子 & 阳 · дома́ \\
парк & 公园 & 阳 · па́рки \\
стадио́н & 体育场 & 阳 · стадио́ны \\
маши́на & 汽车 & 阴 · маши́ны \\
авто́бус & 公共汽车 & 阳 · авто́бусы \\
трамва́й & 有轨电车 & 阳 · трамва́и \\
такси́ & 出租车 & 中 · 不变格 \\
метро́ & 地铁 & 中 · 不变格 \\
ры́нок & 市场 & 阳 · ры́нки \\
пляж & 海滩 & 阳 · пля́жи \\
мо́ре & 海 & 中 · моря́ \\
гора́ & 山 & 阴 · го́ры \\
больни́ца & 医院 & 阴 · больни́цы \\
остано́вка & 车站 & 阴 · остано́вки \\
\hline
\end{longtable}
\begin{longtable}{@{}p{0.27\linewidth}p{0.30\linewidth}p{0.37\linewidth}@{}}
\multicolumn{3}{@{}p{\linewidth}@{}}{\subhdr{3.7　主题词表：大学（课件 p12）}} \\
\hline
\thead{俄语} & \thead{中文} & \thead{性 · 复数} \\
\hline
\endfirsthead
\hline
\thead{俄语} & \thead{中文} & \thead{性 · 复数} \\
\multicolumn{3}{@{}l}{\scriptsize\color{inkgray}（接上页）3.7　主题词表：大学（课件 p12）} \\
\hline
\endhead
факульте́т & 系 & 阳 · факульте́ты \\
библиоте́ка & 图书馆 & 阴 · библиоте́ки \\
аудито́рия & 教室 & 阴 · аудито́рии \\
буфе́т & 小卖部 & 阳 · буфе́ты \\
спортза́л & 体育馆 & 阳 · спортза́лы \\
уро́к & 课 & 阳 · уро́ки \\
ле́кция & 讲座 & 阴 · ле́кции \\
заня́тие & 课、 занятия 课程 & 中 · заня́тия \\
дека́н & 系主任 & 阳 · дека́ны \\
преподава́тель & 教师 & 阳 · преподава́тели \\
гру́ппа & 班级 & 阴 · гру́ппы \\
студе́нт & 大学生 & 阳 · студе́нты \\
студе́нтка & 女大学生 & 阴 · студе́нтки \\
\hline
\end{longtable}
\begin{longtable}{@{}p{0.27\linewidth}p{0.30\linewidth}p{0.37\linewidth}@{}}
\multicolumn{3}{@{}p{\linewidth}@{}}{\subhdr{3.8　主题词表：课堂（课件 p12）}} \\
\hline
\thead{俄语} & \thead{中文} & \thead{性 · 复数} \\
\hline
\endfirsthead
\hline
\thead{俄语} & \thead{中文} & \thead{性 · 复数} \\
\multicolumn{3}{@{}l}{\scriptsize\color{inkgray}（接上页）3.8　主题词表：课堂（课件 p12）} \\
\hline
\endhead
стол & 桌子 & 阳 · столы́ \\
стул & 椅子 & 阳 · сту́лья \\
доска́ & 黑板 & 阴 · до́ски \\
табли́ца & 表格 & 阴 · табли́цы \\
слова́рь & 词典 & 阳 · словари́ \\
уче́бник & 教科书 & 阳 · уче́бники \\
тетра́дь & 练习本 & 阴 · тетра́ди \\
ру́чка & 笔 & 阴 · ру́чки \\
каранда́ш & 铅笔 & 阳 · карандаши́ \\
упражне́ние & 练习 & 中 · упражне́ния \\
кни́га & 书 & 阴 · кни́ги \\
текст & 课文 & 阳 · те́ксты \\
пра́вило & 规则 & 中 · пра́вила \\
зада́ние & 任务、作业 & 中 · зада́ния \\
переры́в & 课间休息 & 阳 · переры́вы \\
\hline
\end{longtable}
\begin{longtable}{@{}p{0.27\linewidth}p{0.30\linewidth}p{0.37\linewidth}@{}}
\multicolumn{3}{@{}p{\linewidth}@{}}{\subhdr{3.9　主题词表：宿舍（课件 p13）}} \\
\hline
\thead{俄语} & \thead{中文} & \thead{性 · 复数} \\
\hline
\endfirsthead
\hline
\thead{俄语} & \thead{中文} & \thead{性 · 复数} \\
\multicolumn{3}{@{}l}{\scriptsize\color{inkgray}（接上页）3.9　主题词表：宿舍（课件 p13）} \\
\hline
\endhead
эта́ж & 楼层 & 阳 · этажи́ \\
ку́хня & 厨房 & 阴 · ку́хни \\
холоди́льник & 冰箱 & 阳 · холоди́льники \\
кастрю́ля & 锅 & 阴 · кастрю́ли \\
сковорода́ & 平底锅 & 阴 · сковороды́ \\
ча́йник & 水壶 & 阳 · ча́йники \\
таре́лка & 盘子 & 阴 · таре́лки \\
нож & 刀 & 阳 · ножи́ \\
ко́мната & 房间 & 阴 · ко́мнаты \\
окно́ & 窗户 & 中 · о́кна \\
ла́мпа & 灯 & 阴 · ла́мпы \\
шкаф & 柜子 & 阳 · шкафы́ \\
по́лка & 架子 & 阴 · по́лки \\
крова́ть & 床 & 阴 · крова́ти \\
компью́тер & 电脑 & 阳 · компью́теры \\
ключ & 钥匙 & 阳 · ключи́ \\
\hline
\end{longtable}
\begin{longtable}{@{}p{0.27\linewidth}p{0.30\linewidth}p{0.37\linewidth}@{}}
\multicolumn{3}{@{}p{\linewidth}@{}}{\subhdr{3.10　主题词表：家庭（课件 p13）}} \\
\hline
\thead{俄语} & \thead{中文} & \thead{性 · 复数} \\
\hline
\endfirsthead
\hline
\thead{俄语} & \thead{中文} & \thead{性 · 复数} \\
\multicolumn{3}{@{}l}{\scriptsize\color{inkgray}（接上页）3.10　主题词表：家庭（课件 p13）} \\
\hline
\endhead
брат & 兄弟 & 阳 · бра́тья \\
сестра́ & 姐妹 & 阴 · сёстры \\
де́душка & 爷爷、外公 & 阳 · де́душки \\
ба́бушка & 奶奶、外婆 & 阴 · ба́бушки \\
дя́дя & 叔叔、舅舅 & 阳 · дя́ди \\
тётя & 阿姨、姑姑 & 阴 · тёти \\
внук & 孙子 & 阳 · вну́ки \\
вну́чка & 孙女 & 阴 · вну́чки \\
ребёнок & 孩子 & 阳 · де́ти \\
\hline
\end{longtable}
\hdr{§4}{УРОК 6：物主代词——按 мой／моя／моё／мои 四栏分}\label{sec:4}
У6 讲 ли́чные и притяжа́тельные местоиме́ния（人称代词与物主代词）。课件的练习全是同一件事：\textbf{Э́то я. Э́то мой го́род, моя́ ма́ма, моё окно́, мои́ роди́тели}——名词的性决定用 мой 还是 моя́。所以这一课的名词不按主题分，\textbf{直接按四栏分}，和课件 p16–p18 的填空顺序一致。这些词绝大多数在 §2、§3 已经出现过，这里只把「配哪个代词」这一层补上；本课真正的新词放在 4.2。
\begin{longtable}{@{}p{0.235\linewidth}p{0.235\linewidth}p{0.235\linewidth}p{0.235\linewidth}@{}}
\multicolumn{4}{@{}p{\linewidth}@{}}{\subhdr{4.1　课件 p16–p18 的名词，按四栏归位}\par{\small 遮住右边三栏，只看一个名词，先判它属于哪一栏，再念出 \textbf{мой／моя́／моё／мои́ ＋ 名词}。课件 p16–p18 的填空就是这条路径；能一口气念对，这一课就过关了。}} \\
\hline
\thead{он → мой} & \thead{она́ → моя́} & \thead{оно́ → моё} & \thead{они́ → мои́} \\
\hline
\endfirsthead
\hline
\thead{он → мой} & \thead{она́ → моя́} & \thead{оно́ → моё} & \thead{они́ → мои́} \\
\multicolumn{4}{@{}l}{\scriptsize\color{inkgray}（接上页）4.1　课件 p16–p18 的名词，按四栏归位} \\
\hline
\endhead
го́род & ма́ма & окно́ & друзья́ \\
университе́т & семья́ & письмо́ & сыновья́ \\
дом & дочь & я́блоко & брю́ки \\
расска́з & дверь & вре́мя & фру́кты \\
слова́рь & жена́ & моро́женое & де́ньги \\
кабине́т & ба́бушка & общежи́тие & рюкзаки́ \\
стол & ко́мната & пальто́ & студе́нты \\
уче́бник & ле́кция & молоко́ & роди́тели \\
каранда́ш & кни́га & ма́сло & о́вощи \\
журна́л & тетра́дь & зада́ние & де́ти \\
гара́ж & ру́чка &  &  \\
уро́к & ла́мпа &  &  \\
де́душка & тётя &  &  \\
день & сестра́ &  &  \\
оте́ц & аудито́рия &  &  \\
дя́дя & газе́та &  &  \\
чай & доска́ &  &  \\
ко́фе & ка́рта &  &  \\
хлеб & маши́на &  &  \\
\hline
\end{longtable}
\begin{longtable}{@{}p{0.27\linewidth}p{0.30\linewidth}p{0.37\linewidth}@{}}
\multicolumn{3}{@{}p{\linewidth}@{}}{\subhdr{4.2　本课新词（前面没出现过的）}} \\
\hline
\thead{俄语} & \thead{中文} & \thead{性 · 复数} \\
\hline
\endfirsthead
\hline
\thead{俄语} & \thead{中文} & \thead{性 · 复数} \\
\multicolumn{3}{@{}l}{\scriptsize\color{inkgray}（接上页）4.2　本课新词（前面没出现过的）} \\
\hline
\endhead
кабине́т & 办公室、诊室 & 阳 · кабине́ты \\
журна́л & 杂志 & 阳 · журна́лы \\
рюкза́к & 背包 & 阳 · рюкзаки́ \\
ма́сло & 黄油、油 & 中 · 常用单数 \\
о́вощи & 蔬菜 & 复 · о́вощи \\
соль & 盐 & 阴 · 常用单数 \\
шокола́д & 巧克力 & 阳 · 常用单数 \\
моро́женое & 冰淇淋 & 中 · 常用单数 \\
\hline
\end{longtable}
\hdr{§5}{УРОК 7：形容词与名词——性、数必须一致}\label{sec:5}
У7 前半（p2–p25）是前六讲复习，p26 起才是新内容：形容词。形容词的词尾跟着名词的性和数走——како́й дом、кака́я у́лица、како́е фо́то、каки́е студе́нты。所以这一课要记的不是「形容词表」，而是\textbf{名词的性}：名词的性一错，后面的形容词一定错。本课新名词按课件的四组练习分。
\begin{longtable}{@{}p{0.27\linewidth}p{0.30\linewidth}p{0.37\linewidth}@{}}
\multicolumn{3}{@{}p{\linewidth}@{}}{\subhdr{5.1　颜色与衣物（课件 p36–p39）}\par{\small 课件在这一组里给的是「颜色形容词 ＋ 名词」的搭配（кори́чневый шарф、кори́чневая су́мка、кори́чневое пальто́、кори́чневые боти́нки），名词只需要记性和数。}} \\
\hline
\thead{俄语} & \thead{中文} & \thead{性 · 复数} \\
\hline
\endfirsthead
\hline
\thead{俄语} & \thead{中文} & \thead{性 · 复数} \\
\multicolumn{3}{@{}l}{\scriptsize\color{inkgray}（接上页）5.1　颜色与衣物（课件 p36–p39）} \\
\hline
\endhead
цвет & 颜色 & 阳 · цвета́ \\
шарф & 围巾 & 阳 · ша́рфы \\
боти́нки & 皮鞋、短靴 & 复 · 单数 боти́нок \\
плащ & 雨衣 & 阳 · плащи́ \\
\hline
\end{longtable}
\begin{longtable}{@{}p{0.27\linewidth}p{0.30\linewidth}p{0.37\linewidth}@{}}
\multicolumn{3}{@{}p{\linewidth}@{}}{\subhdr{5.2　味道与食物（课件 p46）}\par{\small 课件用「甜／酸／苦／咸／辣」五味形容词去问 лимо́н、торт、гру́ша 这些词，名词只要记住性和数。}} \\
\hline
\thead{俄语} & \thead{中文} & \thead{性 · 复数} \\
\hline
\endfirsthead
\hline
\thead{俄语} & \thead{中文} & \thead{性 · 复数} \\
\multicolumn{3}{@{}l}{\scriptsize\color{inkgray}（接上页）5.2　味道与食物（课件 p46）} \\
\hline
\endhead
лимо́н & 柠檬 & 阳 · лимо́ны \\
гру́ша & 梨 & 阴 · гру́ши \\
проду́кты & 食品 & 复 · 单数 проду́кт \\
блю́до & 菜、一道菜 & 中 · блю́да \\
\hline
\end{longtable}
\begin{longtable}{@{}p{0.27\linewidth}p{0.30\linewidth}p{0.37\linewidth}@{}}
\multicolumn{3}{@{}p{\linewidth}@{}}{\subhdr{5.3　人与城市（课件 p34、p40–p42、p47）}\par{\small 这一组是形容词搭配里「人」和「地方」的名词。注意 челове́к 的复数是 лю́ди（§3.2），де́рево 的复数是 дере́вья。}} \\
\hline
\thead{俄语} & \thead{中文} & \thead{性 · 复数} \\
\hline
\endfirsthead
\hline
\thead{俄语} & \thead{中文} & \thead{性 · 复数} \\
\multicolumn{3}{@{}l}{\scriptsize\color{inkgray}（接上页）5.3　人与城市（课件 p34、p40–p42、p47）} \\
\hline
\endhead
челове́к & 人 & 阳 · лю́ди \\
журнали́стка & 女记者 & 阴 · журнали́стки \\
де́вушка & 姑娘 & 阴 · де́вушки \\
лю́ди & 人们 & 复 · 单数 челове́к \\
де́рево & 树 & 中 · дере́вья \\
плечо́ & 肩膀 & 中 · пле́чи \\
кино́ & 电影、电影院 & 中 · 不变格 \\
фильм & 电影 & 阳 · фи́льмы \\
исто́рия & 历史 & 阴 · исто́рии \\
язы́к & 语言 & 阳 · языки́ \\
\hline
\end{longtable}
\begin{longtable}{@{}p{0.27\linewidth}p{0.30\linewidth}p{0.37\linewidth}@{}}
\multicolumn{3}{@{}p{\linewidth}@{}}{\subhdr{5.4　时间与天气（课件 p40、p42、p48）}\par{\small 课件 p48 的四个句子 Мне нра́вится фи́зика／Ей нра́вится спорт／Им нра́вится зима́／Мне нра́вятся фильмы 里，主语都是名词——这一组同时为 §6 的 нра́виться 做铺垫。}} \\
\hline
\thead{俄语} & \thead{中文} & \thead{性 · 复数} \\
\hline
\endfirsthead
\hline
\thead{俄语} & \thead{中文} & \thead{性 · 复数} \\
\multicolumn{3}{@{}l}{\scriptsize\color{inkgray}（接上页）5.4　时间与天气（课件 p40、p42、p48）} \\
\hline
\endhead
у́тро & 早晨 & 中 · 常用单数 \\
год & 年 & 阳 · го́ды \\
ме́сяц & 月份 & 阳 · ме́сяцы \\
ве́тер & 风 & 阳 · ве́тры \\
не́бо & 天空 & 中 · 常用单数 \\
зима́ & 冬天 & 阴 · зи́мы \\
ле́то & 夏天 & 中 · 常用单数 \\
\hline
\end{longtable}
\begin{longtable}{@{}p{0.27\linewidth}p{0.30\linewidth}p{0.37\linewidth}@{}}
\multicolumn{3}{@{}p{\linewidth}@{}}{\subhdr{5.5　学习与生活（课件 p40–p42、p51）}} \\
\hline
\thead{俄语} & \thead{中文} & \thead{性 · 复数} \\
\hline
\endfirsthead
\hline
\thead{俄语} & \thead{中文} & \thead{性 · 复数} \\
\multicolumn{3}{@{}l}{\scriptsize\color{inkgray}（接上页）5.5　学习与生活（课件 p40–p42、p51）} \\
\hline
\endhead
рабо́та & 工作 & 阴 · рабо́ты \\
встре́ча & 见面 & 阴 · встре́чи \\
настрое́ние & 心情 & 中 · настрое́ния \\
спорт & 体育 & 阳 · 常用单数 \\
ша́хматы & 国际象棋 & 复 · ша́хматы \\
фи́зика & 物理 & 阴 · 常用单数 \\
о́пера & 歌剧 & 阴 · о́перы \\
\hline
\end{longtable}
\hdr{§6}{УРОК 8：指示代词与 нра́виться}\label{sec:6}
У8 把「性」又用了一次，这回配的是\textbf{指示代词}：э́тот дом、э́та кни́га、э́то пальто́、э́ти часы́。后半课的 нра́виться 则要求「谁喜欢什么」——喜欢的东西用第一格做主语，动词按主语的数变（нра́вится ／ нра́вятся）。所以这一课的名词仍然按四栏分，和 4.1 同一套思路，只是把代词换成 этот／эта／это／эти。
\begin{longtable}{@{}p{0.235\linewidth}p{0.235\linewidth}p{0.235\linewidth}p{0.235\linewidth}@{}}
\multicolumn{4}{@{}p{\linewidth}@{}}{\subhdr{6.1　课件 p5–p12、p19 的名词，按四栏归位}\par{\small 和 4.1 一样：遮住右边三栏，只看名词，先判四栏，再念 \textbf{э́тот／э́та／э́то／э́ти ＋ 名词}。考试第一大题（Задание 2）就是这四个词的选择题，选项只差一个字母。}} \\
\hline
\thead{э́тот（он）} & \thead{э́та（она́）} & \thead{э́то（оно́）} & \thead{э́ти（они́）} \\
\hline
\endfirsthead
\hline
\thead{э́тот（он）} & \thead{э́та（она́）} & \thead{э́то（оно́）} & \thead{э́ти（они́）} \\
\multicolumn{4}{@{}l}{\scriptsize\color{inkgray}（接上页）6.1　课件 p5–p12、p19 的名词，按四栏归位} \\
\hline
\endhead
телефо́н & арти́стка & пальто́ & часы́ \\
сувени́р & матрёшка & моро́женое & кроссо́вки \\
костю́м & ку́ртка & кафе́ & сувени́ры \\
слова́рь & су́мка & метро́ & цветы́ \\
кот & у́лица & зда́ние & ли́стья \\
фильм & ро́за & не́бо & глаза́ \\
парк & столо́вая & кино́ & де́ньги \\
проспе́кт & пе́сня &  &  \\
го́род & гости́ница &  &  \\
каранда́ш & пого́да &  &  \\
 & о́сень &  &  \\
 & де́вушка &  &  \\
 & маши́на &  &  \\
\hline
\end{longtable}
\begin{longtable}{@{}p{0.27\linewidth}p{0.30\linewidth}p{0.37\linewidth}@{}}
\multicolumn{3}{@{}p{\linewidth}@{}}{\subhdr{6.2　本课新词}} \\
\hline
\thead{俄语} & \thead{中文} & \thead{性 · 复数} \\
\hline
\endfirsthead
\hline
\thead{俄语} & \thead{中文} & \thead{性 · 复数} \\
\multicolumn{3}{@{}l}{\scriptsize\color{inkgray}（接上页）6.2　本课新词} \\
\hline
\endhead
арти́стка & 女演员 & 阴 · арти́стки \\
матрёшка & 套娃 & 阴 · матрёшки \\
костю́м & 西装 & 阳 · костю́мы \\
кроссо́вки & 运动鞋 & 复 · 单数 кроссо́вка \\
ро́за & 玫瑰 & 阴 · ро́зы \\
столо́вая & 食堂 & 阴 · столо́вые \\
поэ́т & 诗人 & 阳 · поэ́ты \\
цвето́к & 花 & 阳 · цветы́ \\
\hline
\end{longtable}
\needspace{6\baselineskip}\subhdr{6.3　нра́виться 的主语：喜欢的东西是「主语」}

课件 p14 的四句是全课的骨架：

\begin{minipage}{\linewidth}
\tabtop
\begin{tabular}{@{}p{0.46\linewidth}p{0.46\linewidth}@{}}
\thead{主语是单数 → нра́вится} & \thead{主语是复数 → нра́вятся} \\
\end{tabular}
\tabmid
\begin{tabular}{@{}p{0.46\linewidth}p{0.46\linewidth}@{}}
Мне нра́вится фи́зика. & Мне нра́вятся кита́йские фи́льмы. \\
Ей нра́вится спорт. & Им нра́вятся э́ти сувени́ры. \\
Им нра́вится зима́. & Ей нра́вятся э́ти цветы́. \\
\end{tabular}
\tabbot
\end{minipage}

\textbf{判断顺序}：先找「喜欢的东西」那个名词，看它是单数还是复数，再决定动词用 \textbf{-тся} 还是 \textbf{-ятся}。注意主语是复数时，前面的第三格代词不变：Мне нра́вятся…、Ей нра́вятся…。
\hdr{§7}{考前速查：一张纸判性、变数}\label{sec:7}
这一节不添新词，只把前面 400 多条词里反复用到的规则压成两页。考前只看这里，加上自己画过道的词，就够了。
\needspace{6\baselineskip}\subhdr{7.1　给一个名词，三步判性}

\begin{minipage}{\linewidth}
\tabtop
\begin{tabular}{@{}p{0.06\linewidth}p{0.30\linewidth}p{0.58\linewidth}@{}}
\thead{步} & \thead{看什么} & \thead{怎么判} \\
\end{tabular}
\tabmid
\begin{tabular}{@{}p{0.06\linewidth}p{0.30\linewidth}p{0.58\linewidth}@{}}
1 & 词尾是 -а／-я？ & 阴性（она́）。例外只有五个指人的词：па́па、де́душка、дя́дя、мужчи́на、ю́ноша，它们是阳性。 \\
2 & 词尾是 -о／-е／-мя？ & 中性（оно́）。-мя 只有 вре́мя、и́мя 几个，复数 -мена。 \\
3 & 词尾是辅音或 -й？ & 阳性（он）。 \\
4 & 词尾是 -ь？ & 判不出来，查 §2.4／§2.6 两栏；口诀见 7.2。 \\
\end{tabular}
\tabbot
\end{minipage}

\subhdr{7.2　-ь 结尾的四条线索（讲义归纳，不是课件原文）}

\begin{itemize}
\item \textbf{ж／ш／ч／щ ＋ ь → 阴性}：ночь、вещь、дочь。因为阳性名词在这四个字母后面不带 ь（нож、мяч、врач、плащ）。
\item \textbf{-ость → 阴性}：но́вость。
\item \textbf{-тель → 阳性}：преподава́тель、писа́тель。
\item \textbf{-арь → 阳性}：слова́рь、календа́рь。
\end{itemize}

四条都不符的，只能查 §2.4 与 §2.6 那两栏，没有别的办法。

\subhdr{7.3　变复数：八个字}

\begin{minipage}{\linewidth}
\tabtop
\begin{tabular}{@{}p{0.36\linewidth}p{0.26\linewidth}p{0.28\linewidth}@{}}
\thead{词干} & \thead{阳性／阴性} & \thead{中性} \\
\end{tabular}
\tabmid
\begin{tabular}{@{}p{0.36\linewidth}p{0.26\linewidth}p{0.28\linewidth}@{}}
硬词干 & -ы & -а \\
软词干 & -и & -я \\
г／к／х／ж／ш／ч／щ 之后 & -и & — \\
\end{tabular}
\tabbot
\end{minipage}

\textbf{三个必背的例外群}：（1）24 组特殊复数（§3.2）；（2）只有单数 8 个（§3.3）；（3）只有复数 17 个（§3.4）。这三群之外，规则都成立。
\begin{longtable}{@{}p{0.27\linewidth}p{0.30\linewidth}p{0.37\linewidth}@{}}
\multicolumn{3}{@{}p{\linewidth}@{}}{\subhdr{7.4　十个最容易判错性的词}\par{\small 这十个词在课件、作业和模拟卷里都出现过，且都不符合「看词尾」的直觉。考前把它们单独抄一遍。}} \\
\hline
\thead{俄语} & \thead{中文} & \thead{性 · 复数} \\
\hline
\endfirsthead
\hline
\thead{俄语} & \thead{中文} & \thead{性 · 复数} \\
\multicolumn{3}{@{}l}{\scriptsize\color{inkgray}（接上页）7.4　十个最容易判错性的词} \\
\hline
\endhead
ко́фе & 咖啡 & 阳 · 不变格 \\
пальто́ & 大衣 & 中 · 不变格 \\
такси́ & 出租车 & 中 · 不变格 \\
де́душка & 爷爷 & 阳 · де́душки \\
мужчи́на & 男人 & 阳 · мужчи́ны \\
день & 白天 & 阳 · дни \\
ночь & 夜 & 阴 · но́чи \\
ме́бель & 家具 & 阴 · 常用单数 \\
о́бувь & 鞋类 & 阴 · 常用单数 \\
часы́ & 钟表 & 复 · часы́ \\
\hline
\end{longtable}
\needspace{6\baselineskip}\subhdr{7.5　考前最后三句话}

\begin{enumerate}
\item \textbf{名词的性不是背出来的，是用出来的}：每记一个名词，嘴上配一次 э́тот／э́та／э́то，比抄十遍有用。
\item \textbf{复数题先看词干软硬，再看七个字母}：г、к、х、ж、ш、ч、щ 后面写 и，其余照 7.3 的表。
\item \textbf{不确定的，回课件对页码}：每个词条都能在 У1–У8 里找到出处；表里的重音和性都照课件抄，不另造。
\end{enumerate}
\hdr{§8}{补充：У8 之后课件里的新名词}\label{sec:8}
У8 之后还有数词（1–1000、价格、时间）与动词三时态（现在／过去／将来）两组课件。它们的主角不是名词，但练习里出现了下面这些新词。期末复习件 ПОВТОРЕНИЕ 不含语音与数词，所以这一节只作补充，不占复习主线。
\begin{longtable}{@{}p{0.27\linewidth}p{0.30\linewidth}p{0.37\linewidth}@{}}
\multicolumn{3}{@{}p{\linewidth}@{}}{\subhdr{8.1　数词课：价格与数量（课件 p5–p9）}\par{\small 数词课的句型是 \textbf{Ско́лько сто́ит／сто́ят ＋ 名词？}，名词要会和 юа́нь／юа́ня／юане́й 三种形式搭配。}} \\
\hline
\thead{俄语} & \thead{中文} & \thead{性 · 复数} \\
\hline
\endfirsthead
\hline
\thead{俄语} & \thead{中文} & \thead{性 · 复数} \\
\multicolumn{3}{@{}l}{\scriptsize\color{inkgray}（接上页）8.1　数词课：价格与数量（课件 p5–p9）} \\
\hline
\endhead
апельси́н & 橙子 & 阳 · апельси́ны \\
ку́рица & 鸡 & 阴 · ку́рицы \\
яйцо́ & 鸡蛋 & 中 · я́йца \\
тома́т & 西红柿 & 阳 · тома́ты \\
носки́ & 袜子 & 复 · 单数 носо́к \\
футбо́лка & T 恤 & 阴 · футбо́лки \\
мяч & 球 & 阳 · мячи́ \\
ча́шка & 杯子 & 阴 · ча́шки \\
ты́сяча & 一千 & 阴 · ты́сячи \\
\hline
\end{longtable}
\begin{longtable}{@{}p{0.27\linewidth}p{0.30\linewidth}p{0.37\linewidth}@{}}
\multicolumn{3}{@{}p{\linewidth}@{}}{\subhdr{8.2　动词三时态课：反复出现的名词（课件各页）}} \\
\hline
\thead{俄语} & \thead{中文} & \thead{性 · 复数} \\
\hline
\endfirsthead
\hline
\thead{俄语} & \thead{中文} & \thead{性 · 复数} \\
\multicolumn{3}{@{}l}{\scriptsize\color{inkgray}（接上页）8.2　动词三时态课：反复出现的名词（课件各页）} \\
\hline
\endhead
баскетбо́л & 篮球 & 阳 · 常用单数 \\
футбо́л & 足球 & 阳 · 常用单数 \\
рок & 摇滚乐 & 阳 · 常用单数 \\
А́нглия & 英国 & 阴 · 常用单数 \\
Белару́сь & 白俄罗斯 & 阴 · 常用单数 \\
Кита́й & 中国 & 阳 · 常用单数 \\
\hline
\end{longtable}
\needspace{6\baselineskip}\subhdr{8.3　这份表的来源与边界}

\begin{itemize}
\item \textbf{词条来源}：УРОК 1--8 共 243 页 ＋ 数词 25 页 ＋ 动词三时态 76 页，合计 344 页课件。每一节标题后的「课件 pXX」都可以回原 PDF 对；同一课里重复出现的词只在该课第一次出现处展开。
\item \textbf{重音与词形}：照课件原文抄；课件没标重音的常用词按标准重音补。复数形式以课件表格为准，课件没给的按规则补，补出来的都在 §3 的规则范围内。
\item \textbf{中文释义}：取课件语境里的常用义，不堆词典义项；一个词在课件里有多个意思时，只给最靠近课文的那一个。
\item \textbf{考试边界}：期末复习件 ПОВТОРЕНИЕ 不含语音与数词，所以 §1 与 §8 只用于「认词」，不扩大考试范围；真正决定分数的是 §2 的性、§3 的数、§4--§6 的代词与形容词一致。
\end{itemize}

\aside{这份表按「认词」设计：先认形、再认性、最后才认复数。三遍读法的具体安排见 §0.3。}
\end{document}
